\chapter{Conclusiones y Trabajo Futuro}
\label{cap:conclusiones}

\section{Conclusiones globales}

El objetivo de este trabajo era construir un sistema capaz de transformar un partido de fútbol en una narración automática, integrando visión por computador, seguimiento de jugadores, identificación de equipos, detección de acciones, generación de comentarios e incorporación de voz. El resultado final no alcanza todavía el comportamiento de una retransmisión profesional en tiempo real, pero sí demuestra que el pipeline completo es viable y que cada una de sus etapas produce una salida suficientemente coherente como para alimentar a la siguiente.

Una de las conclusiones principales del proyecto es que el problema no se resuelve con un único modelo, sino mediante la coordinación de muchos componentes especializados. La detección localiza jugadores, balón y árbitro; el tracking mantiene identidades a lo largo del tiempo; la asignación de equipos permite diferenciar a los jugadores por color; la alineación introducida en la interfaz ayuda a asociar posiciones y nombres; la detección de acciones transforma trayectorias en eventos narrables; el modelo de lenguaje convierte esos eventos en texto; y el sistema de síntesis de voz convierte finalmente ese texto en audio. La dificultad real aparece precisamente en la acumulación de pequeños errores entre módulos.

Aun así, el sistema construido permite recorrer todas las fases del proceso. Se ha conseguido pasar desde un vídeo de entrada hasta un vídeo procesado con comentarios generados automáticamente, alternando voces de comentaristas y permitiendo tanto una reproducción en directo simulada como un ensamblado posterior más alineado temporalmente. 

Los resultados también muestran que la latencia sigue siendo el principal obstáculo para una narración verdaderamente en tiempo real. La GPU utilizada durante el desarrollo, una NVIDIA RTX 3080, ha sido imprescindible para poder entrenar, ejecutar detección, probar modelos de voz y realizar experimentos de inferencia local. Sin embargo, para un proyecto tan ambicioso, en el que varias fases compiten por los mismos recursos de cómputo, se queda limitada. Una GPU más potente reduciría parte de la latencia local, especialmente en detección, tracking y síntesis basada en modelos locales, pero no eliminaría por completo los retrasos introducidos por servicios externos como ElevenLabs, donde existe una latencia de red y de generación que ronda aproximadamente el segundo incluso usando streaming.

En este sentido, una conclusión práctica es que una retransmisión automática completamente sincronizada puede requerir introducir un pequeño retardo controlado en la visualización del vídeo. En un escenario de emisión real, retrasar ligeramente la imagen permitiría disponer de margen para detectar la acción, generar el comentario y reproducirlo de forma más natural. Esta estrategia es habitual en sistemas donde prima la coherencia audiovisual frente a la inmediatez absoluta.

\section{Consideraciones sobre licencias y recursos externos}

Durante el desarrollo se han utilizado herramientas, modelos y conjuntos de datos de terceros. Aunque el proyecto se ha realizado con finalidad académica, no todos estos recursos tienen las mismas condiciones de uso. Algunos son abiertos y pueden utilizarse sin pago, otros son abiertos pero imponen obligaciones fuertes, y otros dependen de servicios comerciales o de licencias específicas.

En la parte de detección se ha utilizado YOLO mediante Ultralytics. Ultralytics ofrece sus modelos y herramientas bajo un esquema dual. La opción abierta es AGPL-3.0, una licencia que permite usar, modificar y distribuir el software sin pagar, pero exige que los proyectos derivados cumplan también las obligaciones de la AGPL. En la práctica, si se integra Ultralytics YOLO en una aplicación distribuida o en un servicio, habría que liberar el código correspondiente bajo una licencia compatible. Si se quisiera utilizar YOLO de Ultralytics en un producto cerrado, propietario o comercial sin abrir el código, sería necesario adquirir una licencia Enterprise de Ultralytics.

Para el tracking se ha utilizado \textit{supervision}, incluyendo la integración con ByteTrack. Esta librería está publicada bajo licencia MIT, que es una licencia abierta y permisiva. Permite uso académico, personal y comercial, modificación y redistribución, sin necesidad de pago, siempre que se mantengan los avisos de copyright y la licencia original.

El conjunto de vídeos de Kaggle \textit{DFL Bundesliga 460 MP4 Videos in 30Sec. + CSV} está publicado bajo licencia CC BY-SA 4.0. Se trata de una licencia abierta que permite usar, compartir y adaptar el material, incluso con fines comerciales, siempre que se cite adecuadamente la fuente. Además, al ser una licencia \textit{ShareAlike}, cualquier adaptación o material derivado que se redistribuya debe hacerse bajo la misma licencia o una compatible. Por tanto, no requiere pago, pero sí atribución y conservación de la misma filosofía de licencia en derivados redistribuidos.

Para el entrenamiento de detección también se ha utilizado el dataset \textit{FootBall Detection} disponible en Roboflow Universe, proporcionado por un usuario de Roboflow. Este dataset está publicado bajo licencia CC BY 4.0. Esta licencia es abierta y más permisiva que CC BY-SA: permite uso, copia, modificación, distribución y uso comercial sin necesidad de pago, siempre que se cite correctamente al autor, se enlace la licencia y se indiquen los cambios realizados si los hubiera. A diferencia de CC BY-SA, no obliga a publicar las adaptaciones bajo la misma licencia.

También se han utilizado datos sintéticos de SoccerSynth/Spiideo SoccerNet SynLoc. En este caso la licencia es más restrictiva: salvo que el dataset indique otra cosa, Spiideo lo ofrece bajo CC BY-NC-SA 4.0 con términos adicionales. Esto significa que puede utilizarse para investigación y experimentación no comercial, citando la fuente y manteniendo la misma licencia en derivados, pero no puede emplearse libremente en un producto comercial. Además, los términos de Spiideo prohíben redistribuir directamente el dataset a terceros, por lo que otros usuarios deberían obtenerlo desde la fuente oficial. Para un uso comercial habría que solicitar autorización o una licencia específica.

Para la detección de acciones se ha integrado PathCRF. El código de PathCRF está publicado bajo licencia MPL-2.0, una licencia abierta que permite uso, modificación y uso comercial sin pago. Sin embargo, no es tan permisiva como MIT: si se modifican archivos cubiertos por MPL-2.0 y se distribuyen, esos archivos modificados deben mantenerse bajo MPL-2.0. PathCRF se apoya originalmente en datos como Sportec Open DFL Dataset, publicado bajo CC BY 4.0. Este dataset permite uso y adaptación, también comercial, siempre que se cite correctamente la fuente.

En la generación de comentarios se ha utilizado Gemma 4, de Google. Gemma 4 está publicado bajo licencia Apache 2.0, una licencia abierta y permisiva. Permite uso académico y comercial, modificación y redistribución sin pago, siempre que se mantengan los avisos de copyright, licencia y atribución correspondientes. Al tratarse de un modelo generativo, además de la licencia conviene respetar las recomendaciones de uso responsable asociadas al modelo.

En la síntesis de voz se han evaluado varias alternativas. XTTS-v2, de Coqui, está publicado bajo la Coqui Public Model License. Esta licencia permite el uso no comercial del modelo y de sus salidas, pero no es una licencia abierta permisiva. Por tanto, XTTS-v2 resulta válido para investigación, pruebas y este proyecto académico, pero no debería utilizarse directamente en un producto comercial sin una licencia o autorización adecuada.

Qwen3-TTS, de Alibaba/Qwen, aparece publicado en Hugging Face bajo licencia Apache 2.0 para los modelos utilizados, incluyendo \texttt{Qwen/Qwen3-TTS-12Hz-1.7B-VoiceDesign} y los modelos \texttt{Base}. En este caso se trata de una licencia abierta y permisiva, por lo que permite uso académico, modificación y uso comercial sin pago, manteniendo los avisos de licencia correspondientes.

ElevenLabs es un caso diferente, ya que no se utiliza como un modelo descargado y ejecutado localmente, sino como un servicio propietario mediante API. No es una solución abierta ni gratuita en producción. Su uso depende de los términos de servicio de ElevenLabs, del sistema de créditos y del plan contratado. Para una demo puede utilizarse con una cuenta y cuota disponible, pero para un uso continuado o comercial habría que asumir el coste del servicio. Además, las voces creadas con Voice Design no se descargan como modelos locales independientes, sino que se consumen dentro de la plataforma de ElevenLabs.

En resumen, los componentes más permisivos del proyecto son \textit{supervision}/ByteTrack, el dataset de Roboflow bajo CC BY 4.0, Sportec Open DFL Dataset, Gemma 4 y Qwen3-TTS. El dataset de Kaggle bajo CC BY-SA 4.0 también es abierto, pero exige mantener la misma licencia en adaptaciones redistribuidas. PathCRF es abierto bajo MPL-2.0, aunque con obligaciones sobre archivos modificados. Ultralytics YOLO puede utilizarse sin pagar bajo AGPL-3.0 si el proyecto cumple sus requisitos de apertura, pero requiere licencia Enterprise para uso propietario cerrado. SoccerSynth/Spiideo está orientado a investigación no comercial y no permite redistribución directa del dataset. XTTS-v2 queda limitado a uso no comercial bajo CPML. ElevenLabs, finalmente, es un servicio propietario de pago por uso o suscripción.

\section{Limitaciones del sistema}

La principal limitación del proyecto es la latencia acumulada. Cada etapa introduce un pequeño coste temporal y, cuando se combinan detección, tracking, identificación, detección de acciones, generación textual y síntesis de audio, el retraso total puede ser perceptible. El sistema funciona mejor cuando puede procesar el vídeo con un pequeño margen o cuando se ensambla el audio en diferido.

Otra limitación importante es que el sistema trabaja con vídeos ya grabados. Aunque se ha simulado un modo en directo, todavía no se ha conectado una cámara real como fuente de entrada continua. Esto implicaría resolver nuevos problemas de captura, buffering, estabilidad de FPS, pérdida de frames, sincronización audiovisual y despliegue en campo.

También existen limitaciones en la detección de acciones. Actualmente se han considerado acciones suficientes para generar una narración básica, pero el fútbol contiene muchos eventos difíciles de reconocer únicamente a partir de trayectorias: faltas, ventajas, tarjetas, fueras de juego, disputas aéreas, cambios de orientación o acciones tácticas más complejas. Algunas de estas acciones requerirían incorporar pose del árbitro, pose de jugadores, contexto temporal más amplio o modelos entrenados específicamente con más datos anotados.

Además, parte del contexto narrativo se encuentra simulado. El sistema permite generar comentarios con información táctica o de clasificación, pero esa información no procede todavía de una base de datos real de competiciones, equipos y jugadores. Esto limita la profundidad de la narración y hace que el comentario sea más dependiente de la información introducida manualmente en la interfaz.

\section{Trabajo futuro}

Una línea de trabajo futura muy relevante sería construir una base de datos amplia para todo el futbol 11 federado en España y luego escalarlo. Esta base de datos podría incluir equipos, plantillas, jugadores, goleadores, clasificaciones, categorías, calendarios y resultados. De esta forma, al lanzar la interfaz no sería necesario escribir manualmente toda la información contextual: bastaría con seleccionar los equipos y el sistema podría preparar automáticamente datos útiles para la narración.

Otra mejora importante sería conectar el sistema a una cámara real. Esto permitiría pasar de vídeos grabados a una retransmisión automática sobre partidos en directo. Para ello habría que diseñar una arquitectura de captura robusta, con procesamiento por ventanas, control de latencia, sincronización entre vídeo y audio, y mecanismos de recuperación ante caídas o frames perdidos.

También sería necesario seguir refinando la detección de acciones. En particular, sería interesante probar modelos específicos para eventos futbolísticos, ampliar el conjunto de acciones reconocidas y estudiar la incorporación de pose del árbitro para detectar faltas, tarjetas o decisiones arbitrales. La pose de los jugadores también podría ayudar a diferenciar controles, tiros, despejes o disputas.

En la parte de audio, una línea futura sería optimizar el sistema para una latencia más baja manteniendo una calidad vocal alta. La combinación de Qwen VoiceDesign con XTTS ha resultado útil como solución local, mientras que ElevenLabs ofrece una calidad superior para demostraciones. En el futuro podrían aparecer modelos TTS locales más rápidos, motores de inferencia optimizados o una estrategia híbrida donde el vídeo se retrase ligeramente para que la narración parezca completamente sincronizada.

En conjunto, el proyecto demuestra que es posible construir un narrador automático de fútbol combinando módulos de visión, modelos de lenguaje y síntesis de voz. El resultado todavía tiene margen de mejora, especialmente en tiempo real, robustez y riqueza contextual, pero constituye una base funcional sobre la que seguir trabajando. La principal aportación no es solo cada módulo por separado, sino haberlos integrado en un pipeline completo capaz de convertir vídeo futbolístico en una narración automática audible.



